\documentclass[11pt]{article}

\usepackage[margin=1in]{geometry}
\usepackage{amsmath,amssymb,amsthm,mathtools}
\usepackage{hyperref}
\usepackage{enumitem}
\usepackage{microtype}

% ---------- Theorem environments ----------
\newtheorem{theorem}{Theorem}[section]
\newtheorem{proposition}[theorem]{Proposition}
\newtheorem{lemma}[theorem]{Lemma}
\newtheorem{corollary}[theorem]{Corollary}
\theoremstyle{definition}
\newtheorem{definition}[theorem]{Definition}
\theoremstyle{remark}
\newtheorem{remark}[theorem]{Remark}

% ---------- Notation shortcuts ----------
\newcommand{\E}{\mathbb{E}}
\newcommand{\Pbb}{\mathbb{P}}
\newcommand{\Var}{\mathrm{Var}}
\newcommand{\Cov}{\mathrm{Cov}}
\newcommand{\R}{\mathbb{R}}
\newcommand{\N}{\mathbb{N}}
\newcommand{\cF}{\mathcal{F}}
\newcommand{\cG}{\mathcal{G}}
\newcommand{\cB}{\mathcal{B}}

\title{Stochastic Calculus --- Class Notes (Proofs \& Details)}
\date{\today}

\begin{document}
\maketitle
\vspace{-1em}

\tableofcontents

% ============================================================
\section{Random Walk and Scaling Limit (Donsker/Invariance Principle)}

\subsection{Symmetric Random Walk (discrete time)}

\begin{definition}[Symmetric random walk]
Let $(\varepsilon_k)_{k\ge1}$ be i.i.d.\ with
\[
\Pbb(\varepsilon_k=1)=\Pbb(\varepsilon_k=-1)=\tfrac12.
\]
Define
\[
X_0=0,\qquad X_n=\sum_{k=1}^n \varepsilon_k,\quad n\ge1.
\]
\end{definition}

\subsubsection*{Slide 23: Draw random walk}
(For the lecture: draw a few sample paths for $n=20,50,200$; emphasize jaggedness.)

\subsubsection*{Slide 24: Mean and variance (proof)}

\begin{proposition}[Mean and variance]
For all $n\ge0$,
\[
\E[X_n]=0,
\qquad
\Var(X_n)=n.
\]
\end{proposition}

\begin{proof}
By linearity,
\[
\E[X_n]=\sum_{k=1}^n \E[\varepsilon_k]=n\cdot 0=0.
\]
Also, the $\varepsilon_k$ are independent with $\Var(\varepsilon_k)=1$, hence
\[
\Var(X_n)=\Var\Big(\sum_{k=1}^n \varepsilon_k\Big)
=\sum_{k=1}^n \Var(\varepsilon_k)
=\sum_{k=1}^n 1
=n.
\]
\end{proof}

\subsection{Continuous-time interpolation and rescaling}

\subsubsection*{Slide 25: Define $W^{(n)}(t)$ (interpolation) and compute $\Var(W^{(n)}(t))$}

\begin{definition}[Scaled random walk with interpolation]
Fix $n\in\N$. For $t\ge0$, define
\[
W^{(n)}(t) \coloneqq \frac{1}{\sqrt{n}}\,X_{\lfloor nt\rfloor}
\quad\text{and extend to all $t$ by linear interpolation on each interval }
\Big[\frac{k}{n},\frac{k+1}{n}\Big].
\]
(Equivalently: if $nt$ is an integer, use $X_{nt}$; otherwise interpolate linearly between the nearest grid points.)
\end{definition}

\begin{proposition}[Independent increments on the grid; mean/variance]
Let $0=t_0<t_1<\cdots<t_m$ with each $nt_j\in\N$.
Then the increments
\[
W^{(n)}(t_1)-W^{(n)}(t_0),\;
W^{(n)}(t_2)-W^{(n)}(t_1),\;\dots,\;
W^{(n)}(t_m)-W^{(n)}(t_{m-1})
\]
are independent. Moreover for $0\le s\le t$ with $ns,nt\in\N$,
\[
\E[W^{(n)}(t)-W^{(n)}(s)]=0,
\qquad
\Var\!\big(W^{(n)}(t)-W^{(n)}(s)\big)=t-s.
\]
\end{proposition}

\begin{proof}
If $ns,nt\in\N$ then
\[
W^{(n)}(t)-W^{(n)}(s)
=\frac{1}{\sqrt{n}}\big(X_{nt}-X_{ns}\big)
=\frac{1}{\sqrt{n}}\sum_{k=ns+1}^{nt}\varepsilon_k.
\]
Disjoint time intervals correspond to disjoint sets of $\varepsilon_k$, hence independence of increments.
Also the expectation is $0$ since each $\varepsilon_k$ is centered.
For the variance:
\[
\Var\!\big(W^{(n)}(t)-W^{(n)}(s)\big)
=\frac{1}{n}\Var(X_{nt}-X_{ns})
=\frac{1}{n}(nt-ns)=t-s.
\]
\end{proof}

\subsubsection*{Slide 26 (often): Martingale property (proof)}

\begin{proposition}[Martingale property (on the grid)]
Let $\cF^{(n)}_t$ be the filtration generated by $(W^{(n)}(u))_{0\le u\le t}$.
For $0\le s\le t$ with $ns,nt\in\N$,
\[
\E\big[W^{(n)}(t)\mid \cF^{(n)}_s\big]=W^{(n)}(s).
\]
\end{proposition}

\begin{proof}
Write $W^{(n)}(t)=W^{(n)}(s)+\big(W^{(n)}(t)-W^{(n)}(s)\big)$.
The increment depends only on $\varepsilon_{ns+1},\dots,\varepsilon_{nt}$,
independent of $\cF^{(n)}_s$, and has mean $0$.
Hence $\E[W^{(n)}(t)-W^{(n)}(s)\mid \cF^{(n)}_s]=0$, so the claim follows.
\end{proof}

\subsubsection*{Slide 27: Donsker theorem (statement you can show)}

\begin{theorem}[Functional CLT / Donsker (informal-to-formal statement)]
As $n\to\infty$, the process $(W^{(n)}(t))_{t\in[0,1]}$ converges in distribution to a standard Brownian motion $(W_t)_{t\in[0,1]}$
(as a random element in $(C([0,1]),\|\cdot\|_\infty)$).

In particular, for any $m\ge1$, any times $t_1,\dots,t_m\in[0,1]$ and any bounded continuous $\varphi:\R^m\to\R$,
\[
\E\big[\varphi\big(W^{(n)}(t_1),\dots,W^{(n)}(t_m)\big)\big]
\;\longrightarrow\;
\E\big[\varphi\big(W(t_1),\dots,W(t_m)\big)\big].
\]
\end{theorem}

\begin{remark}[Pedagogical message]
This is the continuous-time analogue of the CLT: ``a properly rescaled random walk looks like Brownian motion''.
\end{remark}

% ============================================================
\section{Brownian Motion: Definition, Gaussian Structure, Transformations}

\subsection{Definition and first consequences}

\begin{definition}[Standard Brownian motion]
A process $W=(W_t)_{t\ge0}$ is a (standard) Brownian motion if:
\begin{enumerate}[label=(BM\arabic*)]
\item $W_0=0$ a.s.;
\item $W$ has independent increments:
for $0\le t_0<t_1<\cdots<t_m$, the increments $W_{t_1}-W_{t_0},\dots,W_{t_m}-W_{t_{m-1}}$ are independent;
\item Gaussian increments with the right variance:
for $s<t$, $W_t-W_s\sim \mathcal{N}(0,t-s)$;
\item (Path regularity) $t\mapsto W_t(\omega)$ is continuous a.s.
\end{enumerate}
\end{definition}

\begin{proposition}[Mean and covariance]
For Brownian motion,
\[
\E[W_t]=0,
\qquad
\Cov(W_t,W_s)=\min\{t,s\}.
\]
\end{proposition}

\begin{proof}
From (BM3), $W_t=W_t-W_0\sim\mathcal{N}(0,t)$ so $\E[W_t]=0$ and $\Var(W_t)=t$.
For $s\le t$, write $W_t=W_s+(W_t-W_s)$ with $(W_t-W_s)$ independent of $W_s$ and mean $0$, hence
\[
\Cov(W_t,W_s)=\Cov(W_s,W_s)+\Cov(W_t-W_s,W_s)=\Var(W_s)=s.
\]
\end{proof}

\subsection{Gaussian vectors and why mean+covariance characterize them}

\begin{definition}[Gaussian vector]
A vector $Z=(Z_1,\dots,Z_m)\in\R^m$ is Gaussian if every linear combination $\sum_{i=1}^m a_i Z_i$ is (univariate) Gaussian.
\end{definition}

\begin{proposition}[Characterization by mean and covariance]
If $Z$ is Gaussian, then its law is fully determined by
\[
\mu=\E[Z]\in\R^m,
\qquad
\Sigma=\big(\Cov(Z_i,Z_j)\big)_{1\le i,j\le m}\in\R^{m\times m}.
\]
\end{proposition}

\begin{remark}[Gaussian process viewpoint]
Brownian motion is a Gaussian process: for any $t_1,\dots,t_m$, the vector $(W_{t_1},\dots,W_{t_m})$ is Gaussian.
Hence the finite-dimensional distributions are determined by $\E[W_t]=0$ and the covariance kernel $K(s,t)=\min(s,t)$.
\end{remark}

\subsection{Independent increments (sigma-algebra formulation) and Markov}

\subsubsection*{Slide 29: Definition of independent increments}

\begin{definition}[Independent increments via sigma-algebras]
Let $W$ be a process and define the past sigma-algebra
\[
\cF_t \coloneqq \sigma(W_s:\;0\le s\le t).
\]
We say that $W$ has independent increments if for every $t\ge0$,
\[
\sigma\big(W_{t+h}-W_t:\;h\ge0\big)
\quad\text{is independent of}\quad
\cF_t.
\]
\end{definition}

\subsubsection*{Markov property (weak/strong form, clean statement)}

\begin{proposition}[Markov property (Brownian motion)]
For Brownian motion, for every $t\ge0$ and every bounded measurable $f:\R\to\R$,
\[
\E\big[f(W_{t+h})\mid \cF_t\big] = \E\big[f(W_t + (W_{t+h}-W_t))\mid W_t\big]
=: (P_h f)(W_t),
\]
where $(P_h)_{h\ge0}$ is the Brownian transition semigroup.
\end{proposition}

\begin{proof}[Proof sketch]
By independent increments, $W_{t+h}-W_t$ is independent of $\cF_t$ and $\sim \mathcal{N}(0,h)$.
So conditional on $\cF_t$, $W_{t+h}=W_t + \xi$ with $\xi\sim\mathcal{N}(0,h)$ independent of $\cF_t$.
Thus the conditional expectation depends on the past only through $W_t$.
\end{proof}

\begin{proposition}[Strong Markov (statement)]
Let $\tau$ be a stopping time for $(\cF_t)$. Then
\[
(W_{\tau+t}-W_\tau)_{t\ge0}
\quad\text{is independent of } \cF_\tau
\quad\text{and is a Brownian motion.}
\]
\end{proposition}

\begin{proof}[Proof (standard, via approximation by simple stopping times)]
\textbf{Step 0 (strategy).}
We first prove the result for \emph{simple} stopping times (taking finitely many deterministic values),
then approximate a general stopping time by a decreasing sequence of simple stopping times and pass to the limit
using continuity of Brownian paths.

\medskip
\textbf{Step 1 (simple stopping times).}
Assume $\tau$ is simple: $\tau=\sum_{i=1}^m t_i\,\mathbf 1_{A_i}$, where
$0\le t_1<\cdots<t_m$, $(A_i)_{i=1}^m$ is a partition of $\Omega$, and each $A_i\in \cF_{t_i}$.

Fix times $0\le u_1<\cdots<u_n$ and a bounded measurable function
$\varphi:\R^n\to\R$. Define
\[
F := \varphi\big(\widetilde W_{u_1},\ldots,\widetilde W_{u_n}\big)
=\varphi\big(W_{\tau+u_1}-W_\tau,\ldots,W_{\tau+u_n}-W_\tau\big).
\]
To prove independence from $\cF_\tau$ and identify the law, it suffices to show that for every bounded
$\cF_\tau$-measurable random variable $Z$,
\begin{equation}\label{eq:SM-test}
\E\big[ Z\,F\big] \;=\; \E[Z]\;\E\big[\varphi(W_{u_1},\ldots,W_{u_n})\big].
\end{equation}

Because $\cF_\tau$ is generated by $\{A_i\}$ and $\cF_{t_i}$ on each $A_i$,
it is enough (by linearity and a monotone class argument) to verify \eqref{eq:SM-test} for
$Z=\mathbf 1_A$ where $A\in\cF_\tau$. For simple $\tau$, any such $A$ can be written
$A=\bigsqcup_{i=1}^m (A\cap A_i)$ with $A\cap A_i\in \cF_{t_i}$.

Now compute:
\[
\E\big[\mathbf 1_A F\big]
=\sum_{i=1}^m \E\Big[\mathbf 1_{A\cap A_i}\,
\varphi\big(W_{t_i+u_1}-W_{t_i},\ldots,W_{t_i+u_n}-W_{t_i}\big)\Big].
\]
On the event $A\cap A_i\in\cF_{t_i}$, the increment vector
\[
\big(W_{t_i+u_1}-W_{t_i},\ldots,W_{t_i+u_n}-W_{t_i}\big)
\]
is independent of $\cF_{t_i}$ (independent increments) and has the same distribution as
$(W_{u_1},\ldots,W_{u_n})$ (stationary increments). Therefore
\[
\E\Big[\mathbf 1_{A\cap A_i}\,
\varphi(\cdots)\Big]
=\Pbb(A\cap A_i)\;
\E\big[\varphi(W_{u_1},\ldots,W_{u_n})\big].
\]
Summing over $i$ yields
\[
\E\big[\mathbf 1_A F\big]
=\Pbb(A)\;\E\big[\varphi(W_{u_1},\ldots,W_{u_n})\big],
\]
which is exactly \eqref{eq:SM-test} for indicators, hence for all bounded $Z\in L^\infty(\cF_\tau)$.
This proves that the finite-dimensional distributions of $(\widetilde W_t)_{t\ge0}$
match those of Brownian motion and that they are independent of $\cF_\tau$.
In particular, $\widetilde W$ has independent Gaussian increments with variance $t-s$.

\medskip
\textbf{Step 2 (general stopping times by approximation).}
Let $\tau$ be any stopping time. Define a decreasing dyadic approximation
\[
\tau_n := 2^{-n}\,\big\lceil 2^n \tau \big\rceil
\quad\text{so that}\quad
\tau \le \tau_n \downarrow \tau.
\]
Each $\tau_n$ takes values in the grid $\{k2^{-n}:k\in\N\}$ and is a stopping time, and in fact is simple
(on each bounded time horizon; for the full $\tau$ you can localize with $\tau\wedge N$ if desired).
By Step 1, the strong Markov statement holds for each $\tau_n$.

Fix $0\le u_1<\cdots<u_n$ and bounded continuous $\varphi$.
For each $n$, set $\widetilde W^{(n)}_t:=W_{\tau_n+t}-W_{\tau_n}$.
By continuity of Brownian paths and $\tau_n\downarrow\tau$,
\[
\big(\widetilde W^{(n)}_{u_1},\ldots,\widetilde W^{(n)}_{u_n}\big)
\;\xrightarrow[n\to\infty]{}\;
\big(\widetilde W_{u_1},\ldots,\widetilde W_{u_n}\big)
\quad\text{a.s.}
\]
Hence $\,\varphi(\widetilde W^{(n)}_{u_1},\ldots,\widetilde W^{(n)}_{u_n})
\to \varphi(\widetilde W_{u_1},\ldots,\widetilde W_{u_n})$ a.s.,
and dominated convergence applies.

Also, $\cF_\tau\subseteq \cF_{\tau_n}$ (since $\tau\le\tau_n$), so any bounded
$\cF_\tau$-measurable $Z$ is $\cF_{\tau_n}$-measurable. Therefore, by Step 1 applied to $\tau_n$,
\[
\E\!\left[ Z\,\varphi\big(\widetilde W^{(n)}_{u_1},\ldots,\widetilde W^{(n)}_{u_n}\big)\right]
=\E[Z]\;\E\big[\varphi(W_{u_1},\ldots,W_{u_n})\big]
\quad\text{for all }n.
\]
Letting $n\to\infty$ and using dominated convergence gives
\[
\E\!\left[ Z\,\varphi\big(\widetilde W_{u_1},\ldots,\widetilde W_{u_n}\big)\right]
=\E[Z]\;\E\big[\varphi(W_{u_1},\ldots,W_{u_n})\big].
\]
This identity for all bounded $Z\in L^\infty(\cF_\tau)$ and all bounded continuous $\varphi$
implies that the finite-dimensional distributions of $(\widetilde W_t)$ are those of Brownian motion and
that they are independent of $\cF_\tau$.

\medskip
\textbf{Step 3 (conclude Brownian motion).}
We have shown: (i) $\widetilde W_0=0$; (ii) Gaussian stationary independent increments;
(iii) continuity follows from continuity of $W$.
Thus $(\widetilde W_t)_{t\ge0}$ is a Brownian motion independent of $\cF_\tau$.
\end{proof}


\subsection{Transformations of Brownian motion (proofs)}

% ============================================================
% Transformations of Brownian motion
% ============================================================

\begin{proposition}[Symmetry]
If $W=(W_t)_{t\ge0}$ is a Brownian motion, then $\widehat W=(\widehat W_t)_{t\ge0}$ defined by
$\widehat W_t:=-W_t$ is also a Brownian motion.
\end{proposition}

\begin{proof}
We verify the defining properties.
\begin{itemize}[leftmargin=1.5em]
\item $\widehat W_0=-W_0=0$ a.s.
\item (\emph{Independent increments}) For any $0\le t_0<t_1<\cdots<t_n$, the increments of $\widehat W$ are
\[
\widehat W_{t_k}-\widehat W_{t_{k-1}}
=-(W_{t_k}-W_{t_{k-1}}),\qquad k=1,\dots,n.
\]
Since $(W_{t_k}-W_{t_{k-1}})_{k}$ are independent, applying the measurable map
$(x_1,\dots,x_n)\mapsto (-x_1,\dots,-x_n)$ preserves independence, hence the increments of $\widehat W$ are independent.
\item (\emph{Gaussian increments}) For $s<t$,
\[
\widehat W_t-\widehat W_s=-(W_t-W_s).
\]
Because $W_t-W_s\sim \mathcal N(0,t-s)$ and the law $\mathcal N(0,t-s)$ is symmetric, we still have
$\widehat W_t-\widehat W_s\sim \mathcal N(0,t-s)$.
\item (\emph{Continuity}) $t\mapsto \widehat W_t(\omega)=-W_t(\omega)$ is continuous whenever $t\mapsto W_t(\omega)$ is.
\end{itemize}
Thus $\widehat W$ is a Brownian motion.
\end{proof}


\begin{proposition}[Scaling]
Fix $c>0$ and define $\widetilde W=(\widetilde W_t)_{t\ge0}$ by
\[
\widetilde W_t:=\frac1c\,W_{c^2 t}.
\]
Then $\widetilde W$ is a Brownian motion.
\end{proposition}

\begin{proof}
Again we check the defining properties.
\begin{itemize}[leftmargin=1.5em]
\item $\widetilde W_0=\frac1c W_0=0$ a.s.
\item (\emph{Continuity}) $t\mapsto W_{c^2 t}(\omega)$ is continuous, hence so is $t\mapsto \widetilde W_t(\omega)$.
\item (\emph{Independent increments}) Let $0\le t_0<t_1<\cdots<t_n$.
Then
\[
\widetilde W_{t_k}-\widetilde W_{t_{k-1}}
=\frac1c\Big(W_{c^2 t_k}-W_{c^2 t_{k-1}}\Big),\qquad k=1,\dots,n.
\]
Since $c^2 t_0<\cdots<c^2 t_n$, the increments
$\big(W_{c^2 t_k}-W_{c^2 t_{k-1}}\big)_k$ are independent, and scaling by $\frac1c$ preserves independence.
\item (\emph{Gaussian increments with correct variance}) For $s<t$,
\[
\widetilde W_t-\widetilde W_s=\frac1c\big(W_{c^2 t}-W_{c^2 s}\big).
\]
Now $W_{c^2 t}-W_{c^2 s}\sim \mathcal N\big(0,c^2(t-s)\big)$, hence dividing by $c$ gives
\[
\widetilde W_t-\widetilde W_s\sim \mathcal N\big(0,t-s\big).
\]
\end{itemize}
Therefore $\widetilde W$ is a Brownian motion.
\end{proof}


\begin{proposition}[Time reversal]
Fix $T>0$ and define, for $t\in[0,T]$,
\[
\widetilde W_t:=W_T-W_{T-t}.
\]
Then $(\widetilde W_t)_{t\in[0,T]}$ is a Brownian motion on $[0,T]$.
\end{proposition}

\begin{proof}
We verify the Brownian properties on the finite interval $[0,T]$.
\begin{itemize}[leftmargin=1.5em]
\item $\widetilde W_0=W_T-W_T=0$ a.s.
\item (\emph{Continuity}) $t\mapsto W_{T-t}(\omega)$ is continuous on $[0,T]$, hence $\widetilde W$ has continuous paths.
\item (\emph{Independent increments}) Take $0\le t_0<t_1<\cdots<t_n\le T$.
For each $k$,
\[
\widetilde W_{t_k}-\widetilde W_{t_{k-1}}
=\big(W_T-W_{T-t_k}\big)-\big(W_T-W_{T-t_{k-1}}\big)
= W_{T-t_{k-1}}-W_{T-t_k}.
\]
These are increments of $W$ over the disjoint intervals
\[
[T-t_k,\;T-t_{k-1}],\qquad k=1,\dots,n,
\]
which are disjoint because $T-t_n<\cdots<T-t_0$.
Hence by independent increments of $W$, the family
$\big(\widetilde W_{t_k}-\widetilde W_{t_{k-1}}\big)_{k=1}^n$ is independent.
\item (\emph{Gaussian increments with correct variance}) For $0\le s<t\le T$,
\[
\widetilde W_t-\widetilde W_s
= W_{T-s}-W_{T-t}.
\]
This is a Brownian increment over an interval of length $(T-s)-(T-t)=t-s$,
so it is $\mathcal N(0,t-s)$.
\end{itemize}
Thus $\widetilde W$ is Brownian on $[0,T]$.
\end{proof}


\begin{proposition}[Time inversion]
Define $\widetilde W=(\widetilde W_t)_{t\ge0}$ by
\[
\widetilde W_t:= t\,W_{1/t}\quad (t>0),\qquad \widetilde W_0:=0.
\]
Then $\widetilde W$ is a Brownian motion.
\end{proposition}

\begin{proof}
\textbf{Step 1 (Gaussian and covariance).}
For any $t_1,\dots,t_n>0$, the vector $(W_{1/t_1},\dots,W_{1/t_n})$ is Gaussian.
Applying the linear map $(x_i)\mapsto (t_i x_i)$ shows that
\[
(\widetilde W_{t_1},\dots,\widetilde W_{t_n})
=(t_1 W_{1/t_1},\dots,t_n W_{1/t_n})
\]
is also Gaussian.

Moreover, for $s,t>0$,
\[
\Cov(\widetilde W_s,\widetilde W_t)
= st\,\Cov(W_{1/s},W_{1/t})
= st\,\min\{1/s,1/t\}.
\]
If $s\le t$, then $\min\{1/s,1/t\}=1/t$, hence $\Cov(\widetilde W_s,\widetilde W_t)=st\cdot (1/t)=s$.
By symmetry the same formula holds for $t\le s$, so
\[
\Cov(\widetilde W_s,\widetilde W_t)=\min\{s,t\}.
\]
Also $\E[\widetilde W_t]=t\,\E[W_{1/t}]=0$.

\medskip
\textbf{Step 2 (identify finite-dimensional distributions).}
Since $\widetilde W$ is a centered Gaussian process with covariance kernel $\min(s,t)$,
its finite-dimensional distributions coincide with those of standard Brownian motion.

\medskip
\textbf{Step 3 (continuity).}
From Brownian scaling and/or standard modulus-of-continuity estimates, one can show that
$t\mapsto tW_{1/t}$ admits a continuous extension at $0$ with limit $0$ a.s.
(Indeed, $tW_{1/t}\stackrel{d}{=}\sqrt{t}\,W_1$, so it converges to $0$ in $L^2$ and one can upgrade to a.s.
continuity at $0$; away from $0$ it is continuous as a composition of continuous maps.)

\medskip
Therefore $\widetilde W$ has the same law as Brownian motion, hence is a Brownian motion.
\end{proof}

\begin{remark}
The only slightly delicate point is continuity at $t=0$ for the inverted process.
If you prefer a shorter version in lecture notes: state Steps 1--2 (Gaussian + covariance kernel)
and cite the known theorem: ``a centered Gaussian process is determined by its covariance kernel; with this kernel
it admits a continuous modification'', hence it is (a version of) Brownian motion.
\end{remark}
% ============================================================
\subsection{Reflection principle (tool) and hitting time distribution}

\begin{definition}[First passage time]
For $a>0$, define $\tau_a=\inf\{t\ge0:\; W_t=a\}$.
\end{definition}

\begin{proposition}[Reflection principle consequence (statement)]
For $a>0$ and $t>0$,
\[
\Pbb\Big(\sup_{0\le s\le t}W_s \ge a\Big) = 2\Pbb(W_t\ge a)
= 2\Big(1-\Phi\Big(\frac{a}{\sqrt{t}}\Big)\Big),
\]
and therefore
\[
\Pbb(\tau_a\le t)=\Pbb\Big(\sup_{0\le s\le t}W_s \ge a\Big)
=2\Big(1-\Phi\Big(\frac{a}{\sqrt{t}}\Big)\Big).
\]
\end{proposition}

\subsection{Laplace transform of $\tau_a$ (classic exercise)}

\begin{theorem}[Laplace transform of a hitting time]
For $a\in\R$ and $\lambda>0$,
\[
\E\big[e^{-\lambda \tau_a}\big]=\exp\big(-|a|\sqrt{2\lambda}\big).
\]
\end{theorem}

\begin{proof}[Standard proof sketch (exponential martingale + optional stopping)]
Fix $\theta>0$ and consider the exponential martingale
\[
M_t=\exp\Big(\theta W_t-\tfrac12\theta^2 t\Big).
\]
By optional stopping (justified via localization), $\E[M_{\tau_a\wedge t}]=1$.
On $\{\tau_a\le t\}$, $W_{\tau_a}=a$, hence
\[
1=\E[M_{\tau_a\wedge t}]
\ge \E\big[\exp(\theta a-\tfrac12\theta^2 \tau_a)\mathbf{1}_{\{\tau_a\le t\}}\big].
\]
Let $t\to\infty$ (monotone convergence), obtain
\[
1\ge e^{\theta a}\E\big[e^{-(\theta^2/2)\tau_a}\big].
\]
Choose $\theta=\sqrt{2\lambda}$ and rearrange:
\[
\E[e^{-\lambda\tau_a}]\le e^{-\sqrt{2\lambda}\,a}.
\]
Apply the same argument to $-W$ to match $|a|$, and use standard tightness/duality to conclude equality.
(If you want a fully rigorous equality proof: use harmonic functions / ODE for $u(x)=\E_x[e^{-\lambda\tau_a}]$.)
\end{proof}

% 

% ======================================================
% Detailed proofs (put in notes / appendix, not on slides)
% ======================================================

\begin{lemma}[Basic decomposition of the event $\{M_t\ge a\}$]\label{lem:decomp}
Let $a>0$ and $t>0$. Then the events $\{M_t\ge a\}$ and $\{\tau_a\le t\}$ coincide, and
\[
\{M_t\ge a\}
=
\big(\{M_t\ge a,\,W_t\ge a\}\big)\ \dot\cup\ \big(\{M_t\ge a,\,W_t<a\}\big).
\]
\end{lemma}

\begin{proof}
By continuity of $W$, the path reaches level $a$ on $[0,t]$ if and only if the running maximum on $[0,t]$ is at least $a$,
hence $\{M_t\ge a\}\iff\{\tau_a\le t\}$.
The disjoint union is the partition according to whether $W_t\ge a$ or $W_t<a$.
\end{proof}

\begin{lemma}[Reflection map and its key effect]\label{lem:reflection-map}
Fix $a>0$ and $t>0$. On the set $\{\tau_a\le t\}$ define $W^\star$ by
\[
W_s^\star=
\begin{cases}
W_s, & 0\le s\le \tau_a,\\
2a-W_s, & \tau_a< s\le t.
\end{cases}
\]
Then on $\{\tau_a\le t\}$:
\begin{enumerate}[label=(\roman*),leftmargin=2em]
\item $W^\star$ is continuous, $W^\star_{\tau_a}=a$, and $(W^\star)^\star=W$,
\item $\{W_t<a\}$ is mapped to $\{W_t^\star>a\}$ and vice-versa,
\item more generally, for any $w\le a$,
\[
\{W_t\le w\}\ \text{is mapped to}\ \{W_t^\star \ge 2a-w\}.
\]
\end{enumerate}
\end{lemma}

\begin{proof}
(i) Continuity and $W^\star_{\tau_a}=a$ are immediate from continuity of $W$.
Applying the transform twice cancels the reflection, giving $(W^\star)^\star=W$.
(ii) On $\{\tau_a\le t\}$, we have $W_t^\star=2a-W_t$, hence $W_t<a \iff W_t^\star>a$.
(iii) Similarly, $W_t\le w \iff 2a-W_t \ge 2a-w$, i.e.\ $W_t^\star\ge 2a-w$.
\end{proof}

\begin{lemma}[Wiener measure invariance under reflection up to time $t$]\label{lem:invariance}
Fix $a>0$ and $t>0$. For any bounded measurable functional $F$ of the path on $[0,t]$,
\[
\E\big[\,F((W_s)_{0\le s\le t})\,\mathbf 1_{\{\tau_a\le t\}}\,\big]
=
\E\big[\,F((W^\star_s)_{0\le s\le t})\,\mathbf 1_{\{\tau_a\le t\}}\,\big].
\]
\end{lemma}

\begin{proof}[Proof sketch (what you can put in class notes)]
Condition on $\cF_{\tau_a}$. On $\{\tau_a\le t\}$, the post-$\tau_a$ process
\[
\big(W_{\tau_a+u}-W_{\tau_a}\big)_{0\le u\le t-\tau_a}
\]
is (by strong Markov) a Brownian motion independent of $\cF_{\tau_a}$.
By symmetry, this process has the same law as its negative, hence the post-$\tau_a$ segment
has the same conditional law as its sign-flipped version.
But $W^\star$ is exactly the path obtained by keeping the past unchanged and flipping the sign
of the future increments after $\tau_a$ (with a shift by $2a$ to keep continuity at level $a$).
This yields the stated equality of expectations.
(For a fully formal proof: start with cylinder functionals $F$ and use a monotone class argument.)
\end{proof}

\begin{theorem}[Reflection principle]\label{thm:reflection}
Let $a>0$, $t>0$, and $w\le a$. Then
\[
\Pbb(M_t\ge a,\ W_t\le w)=\Pbb(W_t\ge 2a-w).
\]
In particular,
\[
\Pbb(M_t\ge a)=2\,\Pbb(W_t\ge a).
\]
\end{theorem}

\begin{proof}
Using Lemma \ref{lem:decomp},
\[
\Pbb(M_t\ge a,\ W_t\le w)
=
\Pbb(M_t\ge a,\ W_t\le w,\ \tau_a\le t).
\]
On $\{\tau_a\le t\}$ we may apply the reflection map $W\mapsto W^\star$.
By Lemma \ref{lem:invariance},
\[
\Pbb(M_t\ge a,\ W_t\le w)
=
\Pbb(M_t^\star\ge a,\ W_t^\star\le w,\ \tau_a\le t),
\]
and by construction $M_t^\star=M_t$ and $\tau_a(W^\star)=\tau_a(W)$, while Lemma \ref{lem:reflection-map}(iii) gives
\[
\{W_t^\star\le w\}=\{2a-W_t\le w\}=\{W_t\ge 2a-w\}.
\]
Therefore
\[
\Pbb(M_t\ge a,\ W_t\le w)=\Pbb(\tau_a\le t,\ W_t\ge 2a-w).
\]
But if $w\le a$, then $2a-w\ge a$, and the event $\{W_t\ge 2a-w\}$ \emph{already implies} $\{M_t\ge a\}=\{\tau_a\le t\}$,
hence
\[
\Pbb(\tau_a\le t,\ W_t\ge 2a-w)=\Pbb(W_t\ge 2a-w),
\]
which yields the identity.

Taking $w=a$ gives $\Pbb(M_t\ge a,\ W_t\le a)=\Pbb(W_t\ge a)$ and thus
\[
\Pbb(M_t\ge a)
=
\Pbb(M_t\ge a,\ W_t\ge a)+\Pbb(M_t\ge a,\ W_t<a)
=
\Pbb(W_t\ge a)+\Pbb(W_t\ge a)
=
2\,\Pbb(W_t\ge a),
\]
using Lemma \ref{lem:decomp} and the reflection identity with $w=a$.
\end{proof}

\begin{corollary}[First-passage distribution]\label{cor:firstpass}
For $a>0$ and $t>0$,
\[
\Pbb(\tau_a\le t)=2\Pbb(W_t\ge a)
=
\frac{2}{\sqrt{2\pi t}}\int_a^\infty \exp\!\Big(-\frac{x^2}{2t}\Big)\,dx.
\]
\end{corollary}

\begin{proof}
By continuity, $\{\tau_a\le t\}\iff\{M_t\ge a\}$, and apply Theorem \ref{thm:reflection}.
The Gaussian integral is the tail of $\mathcal N(0,t)$.
\end{proof}

\begin{corollary}[Joint density of $(M_t,W_t)$]\label{cor:jointdensity}
For $m>0$ and $w\le m$,
\[
f_{M_t,W_t}(m,w)
=
\frac{2(2m-w)}{t\sqrt{2\pi t}}
\exp\!\Big(-\frac{(2m-w)^2}{2t}\Big).
\]
\end{corollary}

\begin{proof}
For $m>0$ and $w<m$, by Theorem \ref{thm:reflection},
\[
\Pbb(M_t\ge m,\ W_t\le w)=\Pbb(W_t\ge 2m-w)=\int_{2m-w}^\infty p_t(x)\,dx,
\quad p_t(x)=\frac1{\sqrt{2\pi t}}e^{-x^2/(2t)}.
\]
Differentiate in $w$ (Leibniz rule) to get
\[
\frac{\partial}{\partial w}\Pbb(M_t\ge m,\ W_t\le w)=p_t(2m-w).
\]
Differentiate in $m$ and use $p_t'(x)=-(x/t)p_t(x)$:
\[
-\frac{\partial}{\partial m}\,p_t(2m-w)
=
-2p_t'(2m-w)
=
\frac{2(2m-w)}{t}\,p_t(2m-w).
\]
This is exactly the joint density $f_{M_t,W_t}(m,w)$ on the region $m>0$, $w<m$;
the boundary $w=m$ is negligible for densities.
\end{proof}



\subsection{Pathwise Properties (what to prove vs. what to only sketch)}

\subsubsection{At infinity: $\limsup=+\infty$ and $\liminf=-\infty$}

\begin{theorem}[Infinite oscillations at infinity]
With probability $1$,
\[
\limsup_{t\to\infty} W_t = +\infty,
\qquad
\liminf_{t\to\infty} W_t = -\infty.
\]
\end{theorem}

\begin{proof}[Proof sketch (easy, lecture-friendly)]
Use the independence of increments on integer times:
$(W_n)_{n\in\N}$ is a centered random walk with Gaussian steps.
By symmetry and recurrence-type arguments (or Borel--Cantelli on events $\{W_n\ge a\}$ for increasing $a$),
it hits arbitrarily large positive values infinitely often, and by symmetry also arbitrarily large negative values.
\end{proof}

\subsubsection{Near any time: no derivative, huge local slopes (sketch)}
\begin{theorem}[Nowhere differentiability; infinite local slopes]
With probability $1$, the Brownian path $t\mapsto W_t(\omega)$ is nowhere differentiable.
Moreover, for every fixed $t\ge 0$, with probability $1$,
\[
\limsup_{h\downarrow0}\frac{W_{t+h}-W_t}{h}=+\infty,
\qquad
\liminf_{h\downarrow0}\frac{W_{t+h}-W_t}{h}=-\infty.
\]
\end{theorem}

\begin{proof}
\textbf{Step 1 (a deep input: local LIL).}
We use the \emph{local law of the iterated logarithm} (LIL) for Brownian motion:
for every fixed $t\ge 0$, with probability $1$,
\begin{equation}\label{eq:localLIL}
\limsup_{h\downarrow0}\frac{W_{t+h}-W_t}{\sqrt{2h\log\log(1/h)}}=1,
\qquad
\liminf_{h\downarrow0}\frac{W_{t+h}-W_t}{\sqrt{2h\log\log(1/h)}}=-1.
\end{equation}
(We do not reprove \eqref{eq:localLIL} here; it is a classical theorem.)

\medskip
\textbf{Step 2 (deduce $\limsup=+\infty$ and $\liminf=-\infty$ for the slope).}
Fix $t$. From \eqref{eq:localLIL}, along a sequence $h_n\downarrow0$ we have
\[
W_{t+h_n}-W_t \sim \sqrt{2h_n\log\log(1/h_n)}\quad\text{(up to a factor tending to $1$)}.
\]
Therefore
\[
\frac{W_{t+h_n}-W_t}{h_n}
=
\sqrt{\frac{2\log\log(1/h_n)}{h_n}}\cdot
\frac{W_{t+h_n}-W_t}{\sqrt{2h_n\log\log(1/h_n)}}.
\]
The second factor has $\limsup=1$ and $\liminf=-1$ by \eqref{eq:localLIL}, while the prefactor
\[
\sqrt{\frac{2\log\log(1/h_n)}{h_n}}
\;\xrightarrow[n\to\infty]{}\;+\infty
\qquad\text{since } h_n\downarrow0.
\]
Hence the product has $\limsup=+\infty$ and $\liminf=-\infty$, i.e.
\[
\limsup_{h\downarrow0}\frac{W_{t+h}-W_t}{h}=+\infty,
\qquad
\liminf_{h\downarrow0}\frac{W_{t+h}-W_t}{h}=-\infty,
\]
almost surely.

\medskip
\textbf{Step 3 (nowhere differentiability).}
On the event where the previous statement holds for a given $t$, the right difference quotients
$\frac{W_{t+h}-W_t}{h}$ are unbounded above and below as $h\downarrow0$. In particular,
they cannot converge to a finite limit. Thus the \emph{right derivative} at $t$ does not exist.

Applying the same argument to the time-reversed increments
\[
\frac{W_t-W_{t-h}}{h}
=
\frac{W_{(t-h)+h}-W_{t-h}}{h}
\quad (h\downarrow0),
\]
shows the \emph{left derivative} does not exist either. Therefore $W$ is not differentiable at $t$.

Finally, the local LIL \eqref{eq:localLIL} holds (as stated) for every fixed $t$; in fact one has the stronger
result that it holds simultaneously for all $t$ on an event of probability $1$.
Consequently, with probability $1$, $W$ is nowhere differentiable.
\end{proof}

\begin{remark}[Pedagogical message]
Even without the LIL, the key intuition is: typical increments satisfy $|W_{t+h}-W_t|\approx \sqrt{h}$,
so the slope $\frac{|W_{t+h}-W_t|}{h}$ is typically of order $\frac{1}{\sqrt{h}}\to\infty$.
The LIL makes this intuition rigorous (and quantifies the extra $\sqrt{\log\log(1/h)}$ factor).
\end{remark}

\subsubsection{Hölder continuity (replace the LIL modulus slide)}

\begin{theorem}[Hölder continuity]
Fix $T>0$ and $\alpha\in(0,\tfrac12)$. Then with probability $1$, there exists a finite random constant $C_{\alpha,T}(\omega)$ such that
\[
|W_t(\omega)-W_s(\omega)| \le C_{\alpha,T}(\omega)\,|t-s|^{\alpha},
\qquad \forall s,t\in[0,T].
\]
For $\alpha=\tfrac12$, this fails: Brownian motion is not Hölder-$\tfrac12$.
\end{theorem}

\begin{remark}[Proof status]
This is a standard consequence of Kolmogorov's continuity theorem
(using $\E|W_t-W_s|^p = C_p|t-s|^{p/2}$ and taking $p$ large enough so that $p/2>1$).
In class: state it + show the moment scaling $\E|W_t-W_s|^p\propto |t-s|^{p/2}$.
\end{remark}

% ============================================================
\section{Quadratic Variation: Motivation, Definition, Convergence}

\subsection{Motivation (self-financing and why squares matter)}

\paragraph{Goal.}
We want to model a self-financing strategy with underlying $S$ and holdings $\Delta_t$:
\[
V_t = V_0 + \int_0^t \Delta_s\, dS_s.
\]
If we (toy-model) take $S=W$, the ``usual'' Riemann–Stieltjes integral fails because $W$ has infinite variation and is too rough.
The right calculus is It\^o calculus, where the new object is the quadratic variation:
\[
\sum (W_{t_{k+1}}-W_{t_k})^2 \longrightarrow t.
\]
This is the reason why $(dW_t)^2$ behaves like $dt$ in computations.

\subsection{Definition}

\begin{definition}[Quadratic variation along a partition]
Let $\pi=\{0=t_0<t_1<\cdots<t_n=T\}$ be a partition of $[0,T]$ and set $|\pi|=\max_k(t_{k+1}-t_k)$.
Define
\[
[W]^{\pi}_T \coloneqq \sum_{k=0}^{n-1}\big(W_{t_{k+1}}-W_{t_k}\big)^2.
\]
\end{definition}

\begin{remark}
For a $C^1$ function $x(t)$, the same sum goes to $0$ as $|\pi|\to0$.
For Brownian motion it converges to $T$ (non-trivial!).
\end{remark}

\subsection{Convergence (your slides 52: full $L^2$ proof + robustness)}

\begin{proposition}[$L^2$ convergence of quadratic variation]
For Brownian motion $W$ and any partition $\pi$ of $[0,T]$,
\[
\E\Big(\big([W]^{\pi}_T-T\big)^2\Big)\le 2T|\pi|.
\]
In particular, if $|\pi_n|\to0$, then $[W]^{\pi_n}_T\to T$ in $L^2$, hence in probability.
\end{proposition}

\begin{proof}
Let $\Delta W_k:=W_{t_{k+1}}-W_{t_k}\sim\mathcal{N}(0,\Delta t_k)$, independent, with $\Delta t_k=t_{k+1}-t_k$.

\noindent\textbf{Step 1: expectation.}
\[
\E\big[[W]^{\pi}_T\big]=\sum_k \E[(\Delta W_k)^2]=\sum_k \Delta t_k=T.
\]

\noindent\textbf{Step 2: variance.}
For $Z\sim\mathcal{N}(0,\sigma^2)$, one has $\Var(Z^2)=2\sigma^4$.
Thus,
\[
\Var\big([W]^{\pi}_T\big)
=\sum_k \Var\big((\Delta W_k)^2\big)
=\sum_k 2(\Delta t_k)^2
\le 2|\pi|\sum_k \Delta t_k
=2|\pi|\,T.
\]
Since $\E([W]^{\pi}_T)=T$,
\[
\E\Big(\big([W]^{\pi}_T-T\big)^2\Big)=\Var\big([W]^{\pi}_T\big)\le 2T|\pi|.
\]
If $|\pi_n|\to0$, the RHS goes to $0$, giving $L^2$ convergence (and hence convergence in probability).
\end{proof}

\begin{remark}[Optional refinement (almost sure along dyadics)]
One can strengthen to a.s.\ convergence along specific refining sequences (e.g.\ dyadic partitions) using Borel--Cantelli.
Not necessary for first exposure.
\end{remark}




============================================================
\section*{Checklist to match your slide plan}
\begin{itemize}
\item Slide 23: drawing RW (done: note).
\item Slide 24: proof mean/var (done).
\item Slide 25: compute $\Var(W^{(n)}(t))$ and increments (done).
\item Slide 27: Donsker theorem with FDD statement (done).
\item Slide 29: independent increments (sigma-algebra definition) (done).
\item Slide 30: Gaussian vector definition + mean/cov characterize (done).
\item Slide 31: proofs for transformations (symmetry/scaling/time reversal; inversion as remark) (done).
\item Slides 39--43: prove what is easy, skip LIL; provide intuition bounds (done).
\item Slide 46: drawing (kept as note placeholders).
\item Slide 52: $L^2$ convergence of quadratic variation + robustness (done).
\end{itemize}

\end{document}